\section{Introduction}
\label{sec:introduction}

Nuclear dimension and Jiang--Su absorption describe two forms of
regularity in a nuclear $C^*$-algebra. Nuclear dimension counts the
order-zero pieces needed in finite-dimensional completely positive
approximations; an order-zero map preserves orthogonality of positive
elements. Jiang--Su stability means an isomorphism
$A\cong A\otimes\Z$, where $\Z$ is the algebra introduced by Jiang and Su \cite{JS}. The first
condition gives a bound on approximation complexity, while the second
expresses invariance under tensoring with a fixed infinite-dimensional
algebra. Their relationship connects noncommutative dimension with
the structure and classification of nuclear $C^*$-algebras.

The obstruction to passing from finite dimension to absorption is
already visible in a matrix algebra: it has nuclear dimension zero
but does not absorb $\Z$. The relevant exclusion must apply throughout
the ideal structure. A $C^*$-algebra $A$ is \emph{nowhere scattered}
if no quotient $I/J$, for closed two-sided ideals $J\subseteq I\subseteq A$,
is a nonzero elementary algebra $\mathcal K(H)$. This is equivalent
to the absence of finite-dimensional irreducible representations in
every hereditary subalgebra \cite[Theorem~3.1]{TV}.

\begin{theorem}[Nonsimple nuclear regularity]
\label{thm:main}
\label{intro:regularity}
Let $A$ be a separable nuclear complex $C^*$-algebra with no nonzero
elementary ideal subquotients. The following conditions are equivalent:
\begin{enumerate}
\item $\nucdim(A)<\infty$;
\item $A\cong A\otimes\Z$;
\item $\nucdim(A)\le1$.
\end{enumerate}
\end{theorem}

The implication from finite nuclear dimension to absorption proves
Conjecture (C1) of Robert and Tikuisis \cite{RT}. Together the two
implications establish the nuclear-dimension equivalence in the
nonsimple Toms--Winter regularity question of
Schafhauser--Tikuisis--White
\cite[Section~22, Problem~LXXV]{STW99}.
Here and throughout, unadorned tensor products are spatial. The
maximal tensor products used in the independent splitting theorem
below carry an explicit subscript.

The companion \cite[Corollaries~1.3 and~1.4]{CuntzCompanion} adds
Cuntz-semigroup criteria to the main equivalence.

\begin{corollary}[Regularity and Cuntz-semigroup criteria]
\label{cor:regularity-cu}
Under the hypotheses of Theorem~\ref{thm:main}, let
$\iota_A:A\to A\otimes\Z$ be $\iota_A(a)=a\otimes1_\Z$.
The conditions of that theorem are also equivalent to each of the following:
\begin{enumerate}[label=\textup{(\roman*)}]
\item the canonical map $\Cu(\iota_A):\Cu(A)\to\Cu(A\otimes\Z)$ is
an isomorphism in $\mathbf{Cu}$;
\item there is an abstract isomorphism $\Cu(A)\cong\Cu(A\otimes\Z)$
in $\mathbf{Cu}$.
\end{enumerate}
The abstract isomorphism need not be induced by a $*$-homomorphism.
\end{corollary}

The proof, including the nonunital naturality argument for the canonical
map, follows the proof of Theorem~\ref{thm:main} in
Section~\ref{subsec:cu-consequence}.

The dimension construction has a broader formulation, without a
separate assumption on elementary subquotients.

\begin{theorem}[Sharp bound after Jiang--Su absorption]
\label{intro:main}
For every separable nuclear complex $C^*$-algebra $A_0$,
\[
 \nucdim(A_0\otimes\Z)\le1.
\]
\end{theorem}

The bound is optimal uniformly over this class. Nuclear dimension zero
for a separable algebra implies the AF property
\cite[Remarks~2.2(iii)]{WZNuclear}, whereas $\Z$ is unital and
infinite-dimensional with no projections other than $0$ and $1$.
Thus $\nucdim(\Z)=1$.

The absorption implication uses a separate theorem. A positive element
is \emph{full} if it generates the whole ambient algebra as a closed
two-sided ideal, and a \emph{character} is a nonzero $*$-homomorphism
to $\C$.

\begin{theorem}[Full signs in a finite maximal tensor power]
\label{thm:tensor-intro}
For every nonzero unital complex $C^*$-algebra $D$ with no characters,
there is a finite integer $n\ge1$ such that
$D^{\otimes_{\max}n}$ contains two full orthogonal positive elements.
\end{theorem}

No separability or nuclearity is assumed in this theorem. Its finite
conclusion answers affirmatively Kirchberg and R\o rdam's
Question~6.2 \cite{KR}, which asks for this splitting property in the
infinite maximal tensor power. The passage to the infinite power is
by its canonical unital inclusion, which preserves fullness and
orthogonality.

The companion \cite[Theorems~1.1 and~1.2]{Fixed324} obtains another
finite endpoint under the same hypotheses on $D$: for some $m\ge1$,
there is a unital $*$-homomorphism from the prime dimension-drop algebra
$I(2,3)$ to $D^{\otimes_{\max}m}$. For every free ultrafilter $\omega$ on
$\N$, it also proves that a nonzero unital separable complex $C^*$-algebra
$A$ is $\Z$-stable exactly when its norm central-sequence algebra
$A_\omega\cap A'$ has no characters. The proof below uses only the full-sign
endpoint of Theorem~\ref{thm:tensor-intro}.

\subsection{Background and the two obstacles}

Winter and Zacharias introduced nuclear dimension as a noncommutative
covering dimension and established its basic permanence properties
\cite{WZNuclear}. Their order-zero structure theorem
\cite{WZOrderZero} makes a colour into a homomorphism multiplied by a
commuting positive weight. Toms and Winter's theory of strongly
self-absorbing algebras \cite{TW} explains why absorption can supply
commuting tensor factors without changing the isomorphism class.
In the simple setting, Matui--Sato proved that a unital separable simple
nuclear quasidiagonal algebra with a unique trace has decomposition rank
at most one after tensoring with the universal UHF algebra
\cite[Theorem~4.2]{MS}. Decomposition rank requires the same order-zero
decomposition as nuclear dimension, together with contractivity of the
entire return map; its bounds therefore also bound nuclear dimension.
Sato--White--Winter established the bound three
for separable unital simple nuclear $\Z$-stable algebras with a unique trace
\cite[Theorem~B]{SWW}. The two-coloured classification work of
Bosa--Brown--Sato--Tikuisis--White--Winter obtained nuclear dimension
one for non-AF separable unital simple nuclear $\Z$-stable algebras
whose trace simplex is nonempty and Bauer
\cite[Theorem~F]{BBSTWW}. Castillejos--Evington--Tikuisis--White--Winter
removed the trace-boundary restriction in the unital case
\cite[Theorem~B]{CETWW}, and Castillejos--Evington treated the
nonunital case \cite[Theorem~A]{CE}.

Nonsimple antecedents reduce dimension in specific structural
classes. Tikuisis--Winter proved decomposition rank at most two after
$\Z$-stabilization for algebras locally approximable by hereditary
subalgebras of $C(X,\mathcal K)$, with $X$ compact Hausdorff
\cite[Theorem~4.1]{TiW}.
Elliott--Niu--Santiago--Tikuisis extended the bound two to locally
subhomogeneous $\Z$-stable algebras \cite[Theorem~A]{ENST}.
In the purely infinite direction, Szab\'o proved
$\nucdim(A\otimes\mathcal O_\infty)\le3$ for every separable nuclear
$A$ \cite[Theorem~2.3]{SzaboOInfinity}. Bosa--Gabe--Sims--White
then established nuclear dimension at most one for separable nuclear
$\mathcal O_\infty$-stable algebras \cite[Theorem~A]{BGSW}.
Their homomorphism theorem gives the same bound for a nuclear
$\mathcal O_\infty$-stable homomorphism from a separable exact domain
to an arbitrary target \cite[Theorem~C]{BGSW}.

The nonsimple dimension problem has an additional difficulty: a trace
may be bounded on one ideal and infinite outside it, and stably finite
and purely infinite ideal quotients may coexist. The appropriate trace
cone therefore consists of all lower semicontinuous extended traces,
as developed by Elliott--Robert--Santiago \cite{ERS}. Our estimates are
relative to the support rank of a spectral cut generating the ideal
being tested. This rank measures the trace of the support, rather than
the trace value of the cut itself. When it is finite, a relative estimate
controls the error; rank zero forces the error rank to vanish. When it
is infinite, a separate ideal-containment argument forces the comparison
target to have infinite rank. The construction begins with the convex order-zero approximations of
Hirshberg--Kirchberg--White \cite{HKW}, in the refinement of
Brown--Carri\'on--White with asymptotically order-zero incoming maps
\cite{BCW}. It localizes them inside hereditary
corners, and rounds their convex weights to obtain one matrix-valued
order-zero return map.

The comparison stage uses two distinct forms of regularity. At each
coordinate, exactness and $\Z$-stability turn a uniform extended-rank
gap into Cuntz comparison, using R\o rdam's almost-unperforation theorem
\cite{Rordam} and Elliott--Robert--Santiago's comparison criterion
\cite[Proposition~6.2]{ERS}. In the supported relative commutant of the ultrapower,
we instead use abstract dimension functions. Gabe's order-zero
formulation \cite{Gabe} of the Kirchberg--R\o rdam cone criterion
\cite{KR05} supplies finite compression
sums from pointwise ideal containment. Packing these sums yields comparisons
inside that commutant. Fullness, divisibility in a full corner, and
almost unperforation verify a further hereditary-subalgebra condition.
The unitary-equivalence criterion of
Bosa--Brown--Sato--Tikuisis--White--Winter \cite{BBSTWW}, which adapts
Robert--Santiago's positive-element classification method \cite{RS},
turns the comparisons into two weighted conjugacies.

For the reverse implication, Winter proved that a separable simple
unital nonelementary algebra of finite nuclear dimension absorbs $\Z$
\cite[Corollary~6.3]{WinterAbsorption}. Tikuisis extended this to
separable algebras of finite nuclear dimension with finitely many ideals
and no elementary quotient of an ideal
\cite[Corollary~8.7]{TikuisisFiniteIdeals}. Robert--Tikuisis obtained
further nonsimple cases under restrictions on the primitive ideal space
and subquotients \cite[Corollaries~7.11 and~7.16]{RT}, and reduced absorption
for separable algebras of finite nuclear dimension to the existence of
two full orthogonal positive elements in the central sequence algebra
\cite[Theorem~1.2]{RT}. Winter--Zacharias's finite covering-number bound
\cite[Proposition~4.3]{WZNuclear} rules out characters in that algebra
under the nowhere-scattered hypothesis. The remaining step is precisely
a full splitting construction. Kirchberg and R\o rdam isolated the
character and tensor-splitting questions in \cite{KR}. Our finite tensor
construction supplies the needed splitting by combining a uniform
Gaussian sign estimate with pure-state excision. The latter is the
method of Akemann--Anderson--Pedersen \cite{AAP}; we include its short
proof using Kadison transitivity \cite{Kadison}.

Recent Cuntz-semigroup approaches provide a complementary perspective.
Antoine--Perera--Thiel--Vilalta relate finite nuclear dimension and the
Global Glimm Property to comparison and divisibility
\cite[Theorem~6.5]{APTV}; Perera--Thiel--Vilalta characterize absorption
through the corresponding regularity of central sequence algebras
\cite[Theorem~3.6]{PTV}. Our reverse implication instead uses the
finite covering bound and the full-sign theorem to meet the
Robert--Tikuisis criterion directly.

\subsection{Proof overview}

The two implications are independent until the final assembly. Figure~\ref{fig:proof}
records their interfaces.

For the dimension bound, we first work with a unital separable nuclear
$A$, without absorption. In Section~\ref{sec:model}, quasicentral ideal
cuts give finitely many families of hereditary pieces. The pieces are
orthogonal within each fixed threshold family, and averaging the
family compressions retains the prescribed tests. We fix the pieces
and their comparison bounds by finite sums of ideal generators before
choosing inner approximation errors. This ordering turns norm errors into
bounds relative to the support ranks of those generators, uniformly
over all traces for which the ranks are finite.

Section~\ref{sec:traces} places the matrix model in
$A\otimes\Z_1\otimes\Z_2\otimes\Z_3$. The first $\Z$ factor carries
the model. In the norm ultrapower, a positive contraction $q$ identifies
a part on which the model agrees with the original inclusion. The
second factor supplies arbitrarily many orthogonal copies of a norm-one
positive contraction. Relative trace estimates compare the error part
with this common part, and a spare spectral cut bounds the comparison
witnesses before they pass to the ultrapower.

Section~\ref{sec:commutant} constructs a supporting map and uses
approximate domination to pack the error into the common part. This
controls dimension functions in a supported relative commutant without
assuming that commutant exact. A full corner, its divisibility, and
almost unperforation give the hereditary-subalgebra condition needed
for unitary equivalence.
The constant third factor supplies a full-spectrum positive contraction
$k$. Functional-calculus comparison gives unitary conjugacies for the two
weights $k$ and $1-k$. Their sum gives exactly two colours for a corner
inclusion. In Section~\ref{sec:endgame}, stabilization and hereditary
permanence preserve those two colours, and unitization handles the
nonunital algebra. This proves Theorem~\ref{intro:main}.

For absorption, Section~\ref{sec:gaussian} constructs a real tensor
that takes both signs on every product of sufficiently separated pairs
of unit vectors. A Hamming packing provides exponentially many weakly
correlated Gaussian evaluations. Their failure probability decays fast
enough to control a net of all admissible configurations. In
Section~\ref{sec:tensor}, a finite commutator witness converts the
absence of characters into uniform separation of pure-state expectation
vectors. Excision transfers the tensor evaluations to every irreducible
representation of the maximal tensor product. The positive and negative
parts of the resulting selfadjoint element are therefore both full.

Finally, Section~\ref{sec:central} uses the norm central sequence algebra
$F(A)$, with its annihilator quotient in the nonunital case. Finite
covering excludes characters; a separable characterless subalgebra
containing $1_{F(A)}$ has a finite tensor splitting by
Theorem~\ref{thm:tensor-intro}. Commuting
reindexed copies give a unital homomorphism from that tensor power into
$F(A)$. Fullness survives this map, so Robert--Tikuisis's criterion
gives absorption and completes Theorem~\ref{thm:main}.

\begin{figure}[tbp]
\centering
\small
\begin{tabular}{@{}p{0.46\linewidth}@{\hspace{0.06\linewidth}}p{0.46\linewidth}@{}}
\centering \textbf{Dimension bound} &
\centering \textbf{Absorption}\tabularnewline[6pt]
\centering $A$ unital, separable, nuclear\newline
(later $A=\widetilde{A_0}\otimes\Z$) &
\centering $A$ separable, nowhere scattered,\newline $\nucdim(A)<\infty$\tabularnewline[4pt]
\centering $\Downarrow$ & \centering $\Downarrow$\tabularnewline[4pt]
\centering Localized matrix model;\newline
relative ranks and bounded witnesses &
\centering Finite covering makes $F(A)$ characterless;\newline
choose separable characterless $D\subseteq F(A)$\newline
with $1_D=1_{F(A)}$\tabularnewline[4pt]
\centering $\Downarrow$ & \centering $\Downarrow$\tabularnewline[4pt]
\centering Supported commutant comparison;\newline
hereditary condition and weights $k$, $1-k$ &
\centering Independent splitting theorem:\newline
every nonzero unital characterless $D$ has\newline
full orthogonal signs in some $D^{\otimes_{\max}n}$\tabularnewline[4pt]
\centering $\Downarrow$ & \centering $\Downarrow$\tabularnewline[4pt]
\centering Two-colour corner approximation;\newline
stabilization and hereditary permanence &
\centering Reindex commuting copies to obtain\newline
$D^{\otimes_{\max}n}\longrightarrow F(A)$ unital;\newline
the full signs pass to $F(A)$\tabularnewline[4pt]
\centering $\Downarrow$ & \centering $\Downarrow$\tabularnewline[4pt]
\centering $\nucdim(A_0\otimes\Z)\le1$ &
\centering Robert--Tikuisis criterion:\newline
$A\cong A\otimes\Z$\tabularnewline
\end{tabular}
\caption{The two proof branches. The left branch uses spatial tensor
products, first constructs the unital model, and then passes to
$A_0\otimes\Z$ by hereditary permanence. The splitting theorem in the
right column requires neither separability nor nuclearity of $D$;
its application to $F(A)$ uses the displayed unital map.
The branches meet only in the main equivalence.}
\label{fig:proof}
\end{figure}

